\usepackage[top=25mm,bottom=24mm,left=27mm,right=27mm,headheight=15pt,headsep=8mm,footskip=12mm]{geometry}
\usepackage{fontspec,amsmath,amsthm,unicode-math}
\usepackage[english]{babel}
\usepackage{microtype,xcolor,fancyhdr,titlesec,booktabs,longtable,array,tabularx}
\usepackage{graphicx,tikz,enumitem,float,needspace,makeidx}
\makeindex
\usetikzlibrary{arrows.meta,calc,positioning,decorations.pathreplacing,matrix}
\usepackage[unicode,hidelinks]{hyperref}
\usepackage[nameinlink,noabbrev]{cleveref}
\setmainfont{STIXTwoText-Regular.otf}[BoldFont=STIXTwoText-Bold.otf,ItalicFont=STIXTwoText-Italic.otf,BoldItalicFont=STIXTwoText-BoldItalic.otf]
\setmathfont{STIXTwoMath-Regular.otf}
\hypersetup{pdftitle={Gravitation, torsion, and cosmological dynamics},pdfauthor={Oumar Haidara Fall; Rubén Ramos Balsa},pdfsubject={Arithmetic memory, field equations, and transition from contraction to expansion}}
\newcommand{\Autores}{Oumar Haidara Fall \textperiodcentered\ Rubén Ramos Balsa}
\newcommand{\Titulo}{Gravitation, torsion, and cosmological dynamics}
\newcommand{\TituloSegundo}{Arithmetic memory, field equations,\\and transition from contraction to expansion}
\newcommand{\Subtitulo}{}
\newcommand{\RR}{\mathbb R}\newcommand{\ZZ}{\mathbb Z}\newcommand{\QQ}{\mathbb Q}\newcommand{\CC}{\mathbb C}\newcommand{\FF}{\mathbb F}
\newcommand{\Id}{\operatorname{id}}\newcommand{\tr}{\operatorname{tr}}
\newcommand{\Span}{\operatorname{span}}\newcommand{\Cyl}{\operatorname{Cyl}}
\newcommand{\square}{\mdlgwhtsquare}
\providecommand{\dashrightarrow}{\mathrel{\rightdasharrow}}
\newcommand{\TPK}{\mathrm{TPK}}\newcommand{\APP}{\mathrm{APP}}\newcommand{\TRIT}{\{-1,0,+1\}}
\newcommand{\HMT}{\mathrm{HMT}}\newcommand{\Tr}{\operatorname{Tr}}\newcommand{\Spec}{\operatorname{Spec}}
\newcommand{\Cstar}{C^{*}}\newcommand{\dd}{\mathrm d}\newcommand{\R}{\mathbb R}\newcommand{\Z}{\mathbb Z}\newcommand{\I}{\mathrm I}
\definecolor{ink}{HTML}{183D4A}\definecolor{accent}{HTML}{8C3D2E}
\newtheorem{theorem}{Theorem}[section]
\newtheorem{lemma}[theorem]{Lemma}\newtheorem{proposition}[theorem]{Proposition}\newtheorem{corollary}[theorem]{Corollary}
\theoremstyle{definition}\newtheorem{definition}[theorem]{Definition}\newtheorem{axiom}[theorem]{Axiom}
\newtheorem{principle}[theorem]{Principle}
\theoremstyle{remark}\newtheorem{remark}[theorem]{Remark}
\newenvironment{apunte}[1]{\par\Needspace{6\baselineskip}\begin{quote}\small\raggedright\noindent\textbf{Review note: #1.} }{\end{quote}}
\numberwithin{equation}{section}
\setlength{\parindent}{1em}\setlength{\parskip}{3pt plus 1pt minus .5pt}\setlength{\emergencystretch}{2em}
\renewcommand{\normalsize}{\fontsize{10.5}{13.4}\selectfont}\normalsize
\setlist{nosep,topsep=5pt}
\titleformat{\section}{\fontsize{13}{16}\selectfont\bfseries}{\thesection}{.65em}{}
\titleformat{\subsection}{\fontsize{11.2}{14}\selectfont\bfseries}{\thesubsection}{.65em}{}
\titleformat{\subsubsection}{\normalsize\bfseries}{\thesubsubsection}{.65em}{}
\AddToHook{cmd/subsection/before}{\Needspace{7\baselineskip}}
\AddToHook{cmd/subsubsection/before}{\Needspace{6\baselineskip}}
\setcounter{tocdepth}{2}
\makeatletter
\renewcommand*\l@subsection{\@dottedtocline{2}{1.5em}{2.9em}}
\makeatother
\titlespacing*{\section}{0pt}{17pt}{9pt}\titlespacing*{\subsection}{0pt}{13pt}{6pt}
\pagestyle{fancy}\fancyhf{}
\fancyhead[L]{\fontsize{8}{10}\selectfont Gravitation, torsion, and cosmological dynamics}
\fancyhead[R]{}
\fancyfoot[C]{\thepage}\renewcommand{\headrulewidth}{.2pt}
\clubpenalty=10000\widowpenalty=10000\displaywidowpenalty=10000
\raggedbottom
% Reparto tipográfico de cierres breves, sin modificar las fuentes probatorias.
% Secciones 15.2, 16.5 y 18.4: conservar la unidad expositiva en la página siguiente.
\AddToHook{cmd/subsection/before}{%
  \ifnum\value{section}=15 \ifnum\value{subsection}=1 \clearpage\fi\fi
  \ifnum\value{section}=16 \ifnum\value{subsection}=4 \clearpage\fi\fi
  \ifnum\value{section}=18 \ifnum\value{subsection}=3 \clearpage\fi\fi}
% The sectorial corollary has a local space reservation at its statement,
% independent of the numbering of sections.
